\section{平台源码等价市场装配契约（与规则的实现差异）}

\begin{gapnote}
本章保留当前平台求解模型的事实，\textbf{不是南方区域 2025 年 V1.0 规则模型}。
其全系统备用、B$\theta$/DC 电压网络、报价目标及 DC 储能与规则均有差异。
规则正文见第~\ref{sec:southern-rule}~节，逐项结论见
\srcpath{docs/modules/market/southern\_rules\_comparison.md}。
\end{gapnote}

市场实现集中于 \srcpath{src/market/market\_simulation.cpp}（超过六千行）。本章不把它压缩成
教材 SCUC/SCED 通式；只写已逐段核验的报价、关键变量/约束、LODF 和结算赋值，并明确其余覆盖边界。

\subsection{报价生成}

参与者的能量和 commitment 加价分别截断到 $[0,10]$，容量扣留截断到 $[0,0.95]$。
令
\begin{align}
 P_g^{min}&=\max(0,\fld{pmin\_mw}_g),\\
 P_g^{phys}&=\max(P_g^{min},\fld{pmax\_mw}_g),\\
 P_g^{offer}&=P_g^{min}+(P_g^{phys}-P_g^{min})(1-w_g).
\end{align}
最小出力费用、空载、启停报价分别为
\begin{align}
 C_g^{min}&=(c_{2,g}(P_g^{min})^2+c_{1,g}P_g^{min})(1+\mu_g),\\
 C_g^{nl}&=c_{0,g}(1+\mu_g^{commit}),\\
 C_g^{su/sd}&=C_{g,true}^{su/sd}(1+\mu_g^{commit}).
\end{align}
分段数 $K=\max(1,\fld{segment\_count})$，每段宽度
$w=(P_g^{phys}-P_g^{min})(1-w_g)/K$；第 $k$ 段左端
$l_k=P_g^{min}+kw$，价格为
\begin{equation}
 \pi_{gk}=\left[c_{1,g}+\max(0,c_{2,g})(l_k+l_k+w)\right](1+\mu_g).
\end{equation}
这是真实二次成本的弦斜率，并且负 $c_2$ 被截为零；依据为
\srcpath{src/market/participant\_behavior.cpp:cost\_secant\_slope} 和
\srcpath{src/market/participant\_behavior.cpp:submit\_participant\_offers}。

\subsection{SCUC 变量和报价目标}

\srcpath{build\_market\_commitment\_model} 的变量块顺序严格为：机组出力、commitment、startup、
shutdown、备用、能量分段、AC 相角/流/切负荷/弃注入、DC 电压/流/切负荷/弃注入、VSC 双向量与
方向二元、DC/DC 双向量与方向二元、DC 储能充放/能量/方向二元、可选 N-1 总备用。
偏移赋值见 \srcpath{src/market/market\_simulation.cpp:build\_market\_commitment\_model}。

机组目标系数不是抽象 $C_g(p)$，而是
\begin{align}
 c(u_{gt})&=[\max(0,C_g^{min})+\max(0,C_g^{nl})]\Delta t,\\
 c(y_{gt})&=\max(0,C_g^{su}),\quad
 c(z_{gt})=\max(0,C_g^{sd}),\\
 c(r_{gt})&=\max(0,\pi_g^{reserve})\Delta t,\\
 c(p_{gtk}^{seg})&=\max(0,\pi_{gk})\Delta t.
\end{align}
切负荷和弃注入、DC 储能充放报价也直接写入同一线性目标；见
\srcpath{src/market/market\_simulation.cpp:build\_market\_commitment\_model}。

机组出力由
\begin{align}
 p_{gt}&=P_g^{min}u_{gt}+\sum_kp_{gtk}^{seg},\\
 0\le p_{gtk}^{seg}&\le q_{gk}u_{gt},\\
 p_{gt}+r_{gt}&\le P_g^{offer}u_{gt}
\end{align}
定义。状态转换为
\begin{equation}
 u_{gt}-y_{gt}+z_{gt}-\mathbf 1_{t>0}u_{g,t-1}
 =\begin{cases}u_g^0,&t=0,\\0,&t>0,
 \end{cases}
\end{equation}
并有 $y_{gt}+z_{gt}\le1$。爬坡把 MW/min 乘 $60\Delta t$；首时段与当前
\fld{pg\_mw} 比较。最小启停采用向后窗口 startup/shutdown 和，而不是向前强制窗口：
\begin{align}
 \sum_{\tau=\max(0,t-U+1)}^t y_{g\tau}&\le u_{gt},\\
 u_{gt}+\sum_{\tau=\max(0,t-D+1)}^t z_{g\tau}&\le1.
\end{align}
依据为 \srcpath{src/market/market\_simulation.cpp:add\_le（lambda）}。

备用需求使用 AC 与 DC gross demand 之和：
\begin{equation}
 \sum_gr_{gt}\ge\max(0,\fld{upward\_reserve\_fraction})
 (D_t^{AC}+D_t^{DC}).
\end{equation}
启用 N-1 时定义 $R_t=\sum_gr_{gt}$，并对每台 outage 机组创建
$p_{gt}+r_{gt}\le R_t$，见 \srcpath{src/market/market\_simulation.cpp:add\_le（lambda）}。

\subsection{AC/DC 网络平衡}

AC 母线平衡行以切负荷 $+1$、弃注入 $-1$、本地机组 $+1$、VSC AC-to-DC $-1$、
VSC DC-to-AC $+\eta$、出线 $-1$、入线 $+1$，右端为 net AC demand。
AC 线路流等式为
\begin{equation}
 f_{lt}-\frac{S_B}{x_l\tau_l}\theta_{f,t}
 +\frac{S_B}{x_l\tau_l}\theta_{t,t}
 =-\frac{S_B}{x_l\tau_l}\phi_l,
\end{equation}
其中 $|x_l|<10^{-12}$ 时替换为带原符号的 $10^{-6}$，$|\tau_l|\le10^{-12}$ 时取 1，
$\phi_l$ 为弧度。依据为 \srcpath{src/market/market\_simulation.cpp:add\_le（lambda）}。

DC 母线平衡的 VSC、DC/DC、储能和 DC 线路符号逐项见
\srcpath{src/market/market\_simulation.cpp:add\_le（lambda）}。DC 线路实际采用
\begin{equation}
 f_{mt}=\frac{S_B^{dc}}{r_m}(V_{f,t}^{dc}-V_{t,t}^{dc}),
\end{equation}
当 $|r_m|<10^{-12}$ 时一律替换为正 $10^{-6}$，不保留负号。
VSC 和 DC/DC 方向互斥、硬功率/电压/droop 等式以及 DC 储能状态式对应
\srcpath{src/market/market\_simulation.cpp:add\_le（lambda）}。

\subsection{LODF 和 N-1 筛查}

LODF 使用 $b_l=S_B/(x_l\tau_l)$ 装配删去参考母线的稀疏 Bbus；固定相移不进入灵敏度，
但基态流保留相移并产生 warning。对 outage $k$ 求单位 from-to 转移，得到 PTDF 列 $h$；若
$|1-h_k|<10^{-7}$ 或非有限，则把该 outage 记为 islanding 并跳过。否则
\begin{equation}
 LODF_{lk}=\frac{h_l}{1-h_k},\qquad LODF_{kk}=-1.
\end{equation}
依据为 \srcpath{src/market/market\_simulation.cpp:build\_lodf\_model}。

后果流严格为
\begin{equation}
 f_{l|k,t}=f_{l,t}+LODF_{lk}f_{k,t},
\end{equation}
限值为 $\max(0,\fld{n1\_emergency\_rating\_multiplier})\fld{rate\_a\_mva}_l$；
只有 $|f_{l|k,t}|-limit$ 大于 \fld{n1\_violation\_tolerance\_mw} 才报告违规。
依据为 \srcpath{src/market/market\_simulation.cpp:screen\_n1}。

\subsection{结算赋值}

日前客户能量付款按 AC/DC 已服务负荷乘各自 LMP 和 $\Delta t$；外生注入收入按
$(injection-curtailment)\,LMP\,\Delta t$。客户总付款为能量、备用和 uplift charge 之和；
资源总收入为能量、备用和 uplift revenue 之和；拥塞租为客户能量付款减资源能量收入。
依据为 \srcpath{src/market/market\_simulation.cpp:settle\_market}。

实时机组偏差收入不是把全量实时出力按实时价结算，而是
\begin{equation}
 R_g^{dev}=\sum_t(P_{gt}^{RT}-P_{gt}^{DA})\Delta t\,LMP_{b(g),t}^{RT}.
\end{equation}
两结算收入再加日前能量/备用/uplift 和净辅助服务调整：
\begin{align}
 A_g&=R_g^{perf}-C_g^{nonperf}-C_g^{imbalance},\\
 R_g^{2set}&=R_g^{DA,energy}+R_g^{DA,reserve}+U_g+R_g^{dev}+A_g.
\end{align}
备用交付量是指令乘 performance factor；机组惩罚偏差仅在备用激活指令
$P^{reserve}_g>10^{-12}$ 的机组上计算，此时
\begin{equation}
 I_g=\max\!\big(0,\;|P^{RT}-P^{instruction}|-\rho\max(|P^{instruction}|,1)\big),
\end{equation}
无备用激活指令的机组偏差恒取 $0$（\srcpath{src/market/market\_simulation.cpp:5611-5619}），
其中 $P^{instruction}$ 为调度指令出力。
依据均为 \srcpath{src/market/market\_simulation.cpp:run\_real\_time\_market}。

HHI 使用百分数份额而非 $[0,1]$ 份额，且出力与收入两侧的分子分母均经 $\max(0,\cdot)$
截断（\srcpath{src/market/market\_simulation.cpp:4461-4480}）：
\begin{equation}
 HHI_{output}=\sum_p\Big(100\,\frac{\max(0,E_p)}{\sum_q\max(0,E_q)}\Big)^{\!2},
\end{equation}
总量不超过 $10^{-12}$ 时份额为零；出力非负时该截断为空操作，收入 HHI 同理（负收入先截为零），见
\srcpath{src/market/market\_simulation.cpp:aggregate\_participant\_settlement}。

\subsection{验证与校核}

本章装配契约的回归证据为 Catch2 套件 \srcpath{hacdcpf\_test\_market\_simulation}
（32 用例 / 1074 断言、31 通过 / 1 保留失败为缺陷证据，\srcpath{tests/test\_market\_simulation.cpp}），覆盖日前流水线、
定价对偶提取、结算赋值、实时两结算偏差（含
\texttt{Case9 real-time fixed-commitment market closes two-settlement deviations} 与
\texttt{Hybrid real-time market settles DC load forecast deviations}）及 HHI 市场力度量；
铜板切片的 LMP 与结算恒等式另由独立方程 oracle 复核（第~\ref{sec:market-xref}~节）。
本章未逐式展开的分支（DC storage 日循环、warm start 修复、重复博弈等）以其 Catch2 用例为准，
不在本章给出独立数值证据。

\subsection{覆盖边界}

完整 SCUC 还包含 DC storage 日循环、网络限值动态激活、warm start 修复、跨轮求解器状态复用；
固定组合 pricing 模型、对偶提取、AC 验证、重复博弈也各有独立分支。
它们的权威实现分别位于 \srcpath{src/market/market\_simulation.cpp:build\_market\_commitment\_model}、
\srcpath{src/market/market\_simulation.cpp:extract\_pricing}、
\srcpath{src/market/market\_simulation.cpp:run\_day\_ahead\_market}、
\srcpath{src/market/market\_simulation.cpp:run\_repeated\_market\_game}。本章不以标准 SCED 或博弈通式替代尚未逐式展开的代码。
